Terms are listed alphabetically. The number in brackets is the week
where the term is introduced or treated most fully. Notation follows the
course-wide table in Week 1.

{\def\LTcaptype{none} % do not increment counter
\begin{longtable}[]{@{}
  >{\raggedright\arraybackslash}p{(\linewidth - 2\tabcolsep) * \real{0.5000}}
  >{\raggedright\arraybackslash}p{(\linewidth - 2\tabcolsep) * \real{0.5000}}@{}}
\toprule\noalign{}
\begin{minipage}[b]{\linewidth}\raggedright
Term
\end{minipage} & \begin{minipage}[b]{\linewidth}\raggedright
Definition
\end{minipage} \\
\midrule\noalign{}
\endhead
\bottomrule\noalign{}
\endlastfoot
\textbf{Acceptance rate} {[}11{]} & Proportion of Metropolis proposals
that are accepted. Very high: steps too small; very low: steps too
large. A rough target is 0.2--0.5. \\
\textbf{Aggregation artifact} {[}5{]} & A pattern in averaged data that
no individual shows, e.g., a power-law curve from averaging exponential
learners. \\
\textbf{AIC / AICc} {[}7{]} & Akaike information criterion,
\(2\,\mathrm{NLL} + 2k\); estimates out-of-sample predictive accuracy.
AICc adds a small-sample correction. \\
\textbf{Akaike weight} {[}7{]} & \(w_i \propto \exp(-\Delta_i/2)\):
relative support for model \(i\) as the best predictor among the
candidates. \\
\textbf{Attention weight} (\(w_m\)) {[}14{]} & GCM parameter: the share
of attention given to stimulus dimension \(m\) (weights sum to 1). \\
\textbf{Bayes factor} (\(\mathrm{BF}_{12}\)) {[}10{]} & Ratio of
marginal likelihoods
\(p(\mathbf{y}\mid\mathcal{M}_1)/p(\mathbf{y}\mid\mathcal{M}_2)\): how
much the data change the odds between two models. \\
\textbf{Bayesian data analysis} {[}9{]} & Using Bayes' rule to infer
parameters and models from data; a method of the scientist, applicable
to any model. \\
\textbf{Bayesian model of cognition} {[}15{]} & A hypothesis that the
mind performs (approximately) Bayesian inference; a computational-level
theory. \\
\textbf{Beta-binomial model} {[}9{]} & A binomial likelihood with a Beta
prior; the posterior is again a Beta distribution (conjugacy). \\
\textbf{BIC} {[}7{]} & Bayesian information criterion,
\(2\,\mathrm{NLL} + k\ln N\); approximates
\(-2\ln p(\mathbf{y}\mid\mathcal{M})\) and favors simpler models than
AIC. \\
\textbf{Bias--variance trade-off} {[}6{]} & Expected prediction error =
noise + bias² + variance. Flexible models lower bias but raise
variance. \\
\textbf{Bootstrap (parametric)} {[}5{]} & Simulate many datasets from
the fitted model, refit each one, and use the spread of the estimates as
their uncertainty. \\
\textbf{Boundary separation} (\(a\)) {[}1, 14{]} & Random walk / DDM
parameter: the distance between the response boundaries. Represents
caution. \\
\textbf{Burn-in (warm-up)} {[}11{]} & Initial MCMC iterations discarded
because they still reflect the starting values. \\
\textbf{Centered / non-centered parameterization} {[}13{]} & Two
equivalent ways to write a hierarchical model; the non-centered form
(\(d_j = \mu + \sigma z_j\)) avoids funnel geometry and divergences. \\
\textbf{Complete / no / partial pooling} {[}5, 13{]} & One parameter set
for everyone / independent parameters per person / individual parameters
drawn from an estimated group distribution. \\
\textbf{Confusion matrix (model recovery)} {[}7{]} &
\(P(\text{selected model} \mid \text{generating model})\), estimated by
simulation. \\
\textbf{Conjugate prior} {[}9, 12{]} & A prior that yields a posterior
in the same family (e.g., Beta for binomial data); enables Gibbs
sampling. \\
\textbf{Credible interval / HDI} {[}9{]} & An interval containing a
stated posterior probability (e.g., 95\%); the HDI is the narrowest such
interval. \\
\textbf{Criterion} (\(c\)) {[}2{]} & SDT response-bias parameter: the
criterion location relative to the midpoint between the noise and signal
distributions. Positive values are conservative. \\
\textbf{Cross-validation (K-fold, LOO)} {[}6{]} & Fit on part of the
data and score predictions on the held-out part, rotating through \(K\)
folds; LOO holds out one observation at a time. \\
\textbf{\(d'\) (sensitivity)} {[}2{]} & SDT parameter: the distance
between the signal and noise distributions in SD units. \\
\textbf{Delta rule / Rescorla--Wagner rule} {[}15{]} & Error-driven
update \(Q \leftarrow Q + \alpha(r - Q)\). \\
\textbf{Deterministic core} {[}2{]} & The part of a model that gives its
average prediction, \(f(x;\boldsymbol\theta)\). \\
\textbf{Deviance} (\(G^2\)) {[}4{]} &
\(2(\ell_{\text{sat}} - \hat\ell)\): misfit relative to the saturated
model; approximately \(\chi^2\) with df = cells − \(k\). \\
\textbf{Divergent transition} {[}13{]} & A NUTS trajectory whose
numerical integration broke down; it signals unreliable exploration of
part of the posterior. \\
\textbf{Drift rate} (\(v\)) {[}1, 14{]} & Random walk / DDM parameter:
the average rate of evidence accumulation. Represents evidence
quality. \\
\textbf{Drift-diffusion model (DDM)} {[}14{]} & Continuous evidence
accumulation with noise between two boundaries; predicts choices and RT
distributions. \\
\textbf{Effective sample size (ESS)} {[}12{]} & The number of
independent draws that would give the same precision as the
autocorrelated MCMC draws. \\
\textbf{Error surface} {[}3{]} & The discrepancy (SSE or NLL) as a
function of the parameters; its shape reveals minima, valleys
(trade-offs) and plateaus. \\
\textbf{Evidence / marginal likelihood} {[}9, 10{]} &
\(p(\mathbf{y}\mid\mathcal{M}) = \int p(\mathbf{y}\mid\boldsymbol\theta)p(\boldsymbol\theta)\,d\boldsymbol\theta\):
how well the model predicted the data before fitting. \\
\textbf{EZ-diffusion} {[}14{]} & Closed-form estimates of drift,
boundary and non-decision time from accuracy and the mean and variance
of correct RTs. \\
\textbf{Flexibility (model complexity)} {[}6{]} & The range of data
patterns a model can produce; it depends on the number of parameters
\emph{and} the functional form. \\
\textbf{Forward / inverse problem} {[}1{]} & Simulation (parameters \ensuremath{\rightarrow}
data) vs.~estimation (data \ensuremath{\rightarrow} parameters). \\
\textbf{Full conditional} {[}12{]} & The distribution of one parameter
given all others and the data; Gibbs sampling draws from it. \\
\textbf{GCM (Generalized Context Model)} {[}14{]} & Exemplar model of
categorization: classification by summed similarity to stored
exemplars. \\
\textbf{Generalization criterion} {[}6{]} & Evaluating models by
predicting new conditions (extrapolation) rather than new trials. \\
\textbf{Gibbs sampling} {[}12{]} & MCMC that updates each parameter by
drawing from its full conditional; every draw is accepted. \\
\textbf{Grid approximation} {[}9{]} & Computing the posterior on a grid
of parameter values; feasible for 1--2 parameters. \\
\textbf{Hamiltonian Monte Carlo (HMC) / NUTS} {[}13{]} & Gradient-based
MCMC that simulates physical dynamics to make distant, high-acceptance
proposals; NUTS tunes the trajectory length automatically. \\
\textbf{Hessian / observed Fisher information} {[}5{]} & Matrix of
second derivatives of the NLL at the MLE; its inverse approximates the
covariance of the estimates. \\
\textbf{Hierarchical (multilevel) model} {[}13{]} & A model with
individual-level parameters drawn from group-level distributions. \\
\textbf{Identifiability (structural / practical)} {[}7{]} & Whether
parameters can be uniquely determined, in principle (structural) or from
the available data (practical). \\
\textbf{Ideal observer} {[}15{]} & A model of optimal inference given
the available information, e.g., reliability-weighted cue
combination. \\
\textbf{Inverse temperature} (\(\beta\)) {[}15{]} & Softmax parameter:
the sensitivity of choices to value differences (exploration
vs.~exploitation). \\
\textbf{Inversion matrix (model recovery)} {[}7{]} &
\(P(\text{generating model} \mid \text{selected model})\): how much to
trust a model-selection outcome. \\
\textbf{Learning rate} (\(\alpha\)) {[}15{]} & How strongly prediction
errors update values. \\
\textbf{Likelihood} (\(L\)), \textbf{log-likelihood} (\(\ell\)),
\textbf{NLL} {[}4{]} & \(p(\mathbf{y}\mid\boldsymbol\theta)\) viewed as
a function of \(\boldsymbol\theta\); its log; and the negative log,
which optimizers minimize. \\
\textbf{Likelihood-ratio test (LRT)} {[}7{]} & For nested models,
\(2\Delta\ell \sim \chi^2_{\Delta k}\) under the simpler model. \\
\textbf{Linking assumption} {[}2{]} & How model quantities map onto
observable behavior. \\
\textbf{Local minimum} {[}3{]} & A point lower than its surroundings but
not the global minimum; found by single-start optimization. Use
multi-start. \\
\textbf{MAP} {[}9{]} & Maximum a posteriori estimate (posterior mode);
equals the MLE under a flat prior. \\
\textbf{Markov chain Monte Carlo (MCMC)} {[}11{]} & Sampling from a
posterior by simulating a Markov chain whose stationary distribution is
the posterior. \\
\textbf{Maximum likelihood estimate (MLE)} {[}4{]} & The parameter
values that maximize the likelihood. \\
\textbf{MCSE} {[}12{]} & Monte Carlo standard error, posterior
SD/\(\sqrt{\mathrm{ESS}}\): the simulation error of a posterior mean. \\
\textbf{Metropolis / Metropolis--Hastings} {[}11{]} & MCMC that proposes
a move and accepts it with probability \(\min(1, \text{ratio})\);
Hastings adds a correction for asymmetric proposals. \\
\textbf{Model mimicry} {[}7{]} & One model producing data that look as
if generated by another. \\
\textbf{Model recovery} {[}7{]} & Simulating from each candidate model
and checking whether the comparison method selects the generating
model. \\
\textbf{Monte Carlo integration} {[}11{]} & Approximating an expectation
by the average over random draws; error \ensuremath{\propto} \(1/\sqrt{S}\). \\
\textbf{Multi-start} {[}3{]} & Running an optimizer from many starting
values and keeping the best result. \\
\textbf{Nelder--Mead simplex} {[}3{]} & Derivative-free optimizer that
moves a simplex by reflection, expansion, contraction and shrinkage. \\
\textbf{Nested models} {[}7{]} & One model is a special case of the
other (obtained by fixing parameters). \\
\textbf{Noise model} {[}2{]} & The distribution of observations around
the deterministic prediction (e.g., binomial, Gaussian). \\
\textbf{Non-decision time} (\(t_0\), \(T_{er}\)) {[}1, 14{]} & The part
of RT not spent accumulating evidence (encoding and motor execution). \\
\textbf{Overfitting} {[}6{]} & Fitting the noise of a sample, which
improves fit but worsens prediction of new data. \\
\textbf{Parameter recovery} {[}7{]} & Simulating data from known
parameters at the actual design and checking that fitting returns
them. \\
\textbf{Parameter trade-off} {[}3, 5{]} & Changes in one parameter can
be compensated by another, which shows up as a tilted valley in the
error surface and correlated estimates. \\
\textbf{Posterior} {[}9{]} &
\(p(\boldsymbol\theta\mid\mathbf{y}) \propto p(\mathbf{y}\mid\boldsymbol\theta)p(\boldsymbol\theta)\). \\
\textbf{Posterior predictive check (PPC)} {[}10{]} & Comparing data
simulated from the fitted model with the observed data, to check
adequacy. \\
\textbf{Power / exponential forgetting} {[}2{]} & \(A(1+t)^{-d}\)
(decreasing relative loss rate) vs.~\(Ae^{-\lambda t}\) (constant
relative loss rate). \\
\textbf{Prior / prior predictive check} {[}9{]} & The distribution of
parameters before the data / simulating data from the prior to check
that the prior is sensible. \\
\textbf{Profile likelihood} {[}5{]} & The NLL minimized over the other
parameters for each value of the parameter of interest; gives
likelihood-ratio intervals. \\
\textbf{Prototype model} {[}14{]} & Categorization by similarity to a
single summary representation per category. \\
\textbf{R-hat} (\(\hat R\)) {[}12{]} & Ratio comparing between-chain and
within-chain variance; \ensuremath{\approx} 1 (\textless{} 1.01) when chains agree. \\
\textbf{Random walk model} {[}1{]} & Discrete-step evidence accumulation
to a boundary; the precursor of the DDM. \\
\textbf{Relative loss rate (hazard)} {[}2{]} & \(-f'(t)/f(t)\): the
proportional rate of forgetting at time \(t\). \\
\textbf{Reparameterization} {[}3, 13{]} & Optimizing or sampling on a
transformed scale (e.g., \(\lambda = e^{\eta}\)) to enforce constraints
or improve geometry. \\
\textbf{RMSD} {[}3{]} & Root mean squared deviation,
\(\sqrt{\mathrm{SSE}/N}\), in data units. \\
\textbf{ROC / zROC} {[}2{]} & Hit rate vs.~false-alarm rate as the
criterion varies / the same on \(z\) (probit) axes. \\
\textbf{Saturated model} {[}4{]} & A model with one free parameter per
data cell, which reproduces the data exactly. \\
\textbf{Savage--Dickey density ratio} {[}10{]} & For nested models,
\(\mathrm{BF}_{01}\) = posterior density / prior density at the null
value. \\
\textbf{Shrinkage} {[}13{]} & The pull of individual estimates toward
the group mean in hierarchical models, stronger for noisier data. \\
\textbf{Signal detection theory (SDT)} {[}2{]} & A model of detection
decisions with noisy evidence and a criterion; separates sensitivity
from bias. \\
\textbf{Softmax} {[}15{]} & \(P(o) \propto \exp(\beta Q(o))\): turns
values into choice probabilities. \\
\textbf{SSE / MSE} {[}3{]} & Sum / mean of squared errors between data
and predictions. \\
\textbf{Standard error} {[}5{]} & The SD of an estimator's sampling
distribution. \\
\textbf{Wald interval} {[}5{]} & \(\hat\theta \pm 1.96\,\mathrm{SE}\),
with the SE from the Hessian; assumes approximate normality. \\
\textbf{Weighted least squares} {[}3{]} & Squared errors weighted by the
inverse of their variance. \\
\end{longtable}
}

\subsection{AI \& LLM terms (from the ``Connection to AI \& LLMs''
sections)}

{\def\LTcaptype{none} % do not increment counter
\begin{longtable}[]{@{}
  >{\raggedright\arraybackslash}p{(\linewidth - 2\tabcolsep) * \real{0.5000}}
  >{\raggedright\arraybackslash}p{(\linewidth - 2\tabcolsep) * \real{0.5000}}@{}}
\toprule\noalign{}
\begin{minipage}[b]{\linewidth}\raggedright
Term
\end{minipage} & \begin{minipage}[b]{\linewidth}\raggedright
Definition
\end{minipage} \\
\midrule\noalign{}
\endhead
\bottomrule\noalign{}
\endlastfoot
\textbf{Benchmark contamination} {[}6{]} & Test items (or their answers)
that appeared in a model's training data; a form of data leakage that
inflates scores. \\
\textbf{Bradley--Terry model} {[}15{]} &
\(P(i \succ j) = 1/(1 + e^{-(r_i - r_j)})\): choice model for paired
comparisons; used for RLHF reward models. Equivalent to a two-option
softmax with \(\beta = 1\). \\
\textbf{Cross-entropy loss} {[}4{]} & The training loss of language
models: the negative log-likelihood of the next token, averaged over
tokens. \\
\textbf{Embedding} {[}14{]} & A vector that a model assigns to a word or
text; distances between embeddings are often used as a proxy for
similarity. \\
\textbf{In-context learning} {[}7{]} & An LLM's improvement on a task
from examples given in the prompt, without changing its parameters. \\
\textbf{Large language model (LLM)} {[}1{]} & A neural network (usually
a transformer) trained to predict the next token in very large text
corpora. \\
\textbf{Perplexity} {[}4{]} & \(\exp(\text{NLL per token})\): the
effective number of equally likely next tokens; lower is better. \\
\textbf{RLHF} {[}15{]} & Reinforcement learning from human feedback: a
reward model is fitted to human preference choices, and the language
model is then trained to maximize that reward. \\
\textbf{Scaling law} {[}3{]} & An empirical (usually power-law) relation
between model loss and model size, data size or compute. \\
\textbf{Self-consistency} {[}11{]} & Sampling several answers from a
model and taking the majority vote. \\
\textbf{Temperature} {[}1{]} & Divides the logits before the softmax in
an LLM; low temperature gives near-deterministic output, high
temperature more random output (\(\beta = 1/T\)). \\
\textbf{Token} {[}1{]} & The unit a language model reads and predicts: a
word or part of a word. \\
\end{longtable}
}
